\section{Experimental Provenance and Control Details}
\label{app:experimental-details}
\subsection{Sources of the Reported Aggregates}
\begin{table}[t]
\caption{Earlier GeoToolEval baseline comparison. AgentGS process metrics are per-record averages from the present 603-record end-to-end experiment; FAS values are retained from the source manuscript. Direct-answer language models have no tool traces; their process metrics are not applicable (---).}
\label{tab:prior-results}
\centering
\begin{tabularx}{\columnwidth}{@{}Xrrr@{}}
\toprule
Method & FAS & TSF1 & TOS \\
\midrule
AgentGS (GPT-5.6 Terra) & \textbf{0.8731} & 0.9170 & 0.9166 \\
GeoAgent (GPT-5.6 Terra) & 0.7804 & 0.8663 & 0.7911 \\
EarthAgent (GPT-5.6 Terra) & 0.0862 & 0.1182 & 0.1390 \\
DeepSeek V4 Pro & 0.5162 & --- & --- \\
GPT-5.6 Terra & 0.3137 & --- & --- \\
GPT-5.4 & 0.3069 & --- & --- \\
Qwen-Plus & 0.2918 & --- & --- \\
GeoGPT & 0.2500 & --- & --- \\
\bottomrule
\end{tabularx}
\end{table}

Table~\ref{tab:prior-results} retains the answer scores supplied by the earlier experiment, with AgentGS identified as its first complete 603-record run. AgentGS's TSF1/TOS are supplied by the recorded end-to-end evaluation on that population. Combining these aggregates in one table does not establish paired answer/process outcomes for every task.

The 57-record configuration study has a separate Full reference. Full's process scores are taken from the supplementary experiment, while A1--A6 use per-task records in the formal worksheet. A1's partial control, A2's inactive intervention, and A5's partial ranking change apply to those reported runs. Current implementation of a corresponding switch does not validate the identity or behavior of a past run.

GeoBenchX supplies two 202-task aggregates and uses GPT-5.6 Terra as both backbone and single reference-answer judge. The exact paired task identities and matched execution settings are not established by the aggregate table. Its 0--2 scoring scale must remain separate from GeoToolEval's answer-type-specific FAS.

The available manuscript evidence does not provide a complete set of immutable model snapshots, prompts, sampling settings, per-task budgets, source revisions, and run identifiers for all comparisons. Those values are not inferred from the current runtime defaults. The present tables report available measurements without substituting newly implemented controls for unverified historical conditions.

\subsection{Interventions in the Current Runtime}
Module interventions operate within Skill-guided planning. A1 preserves access to strategy declarations and admitted capabilities while requesting tool-instance planning without an ordered capability plan. This differs from a tool-only baseline that removes the Skill/Capability layer altogether.

A2 disables automatic predecessor closure during preparation and support completion. Support insertion does not restore omitted predecessors. A3 suppresses supporting-capability insertion and bridging, while retaining gap diagnosis. With A2 and A3 together, neither automatic insertion phase supplies missing dependencies. The model may still explicitly include declared requirements. Effective intervention therefore needs to be evaluated through changes in preparation and planning, not assumed from a configuration flag.

A4 suppresses only directed bridging. A5 selects the first eligible implementation in a fixed tool-identifier order, preserving instance identities, constraints, and argument checks. A6 suppresses execution-driven repair and replanning after a scientific failure. These definitions specify intended controlled comparisons; the reported A1, A2, and A5 measurements retain the limitations stated above.

Effective settings, requested interventions, and suppressed controls are recorded with execution provenance. A resume operation must preserve the same experiment identity. These records support checking whether a run applies its requested intervention but do not establish a nonzero effect on a task with no relevant dependency, alternative tool, or recoverable failure.

\subsection{Candidate Access, Budgets, and Additional Baselines}
A comparison of planning representations should specify whether all systems receive the same tool candidates or retain their own retrieval policies. The runtime supports exposing all registered candidates, while the retrieved setting preserves system-specific candidate selection. Routing hints and binding assistance are separate controls and can remain active even when heuristic plan generation is disabled.

A Bash-based data-only baseline changes the execution interface more substantially. A comparable data-access scope excludes materialized-query lookup tasks whose snapshot service is unavailable in that environment. A comparison that uses the full tool population for one system and a restricted accessible population for another must not interpret the difference as a planning effect. No new Bash-baseline score is asserted here.

Budgets should distinguish model requests, transport attempts, plan attempts, scientific executions, and diagnostic calls. The runtime records separate counters, but the supplied result tables do not include matched latency, token usage, or execution-cost measurements. Consequently, they establish no quantitative cost advantage.

\subsection{Trace-Level Analyses Not Established by Aggregates}
To relate composition to outcomes, a trace must show the selected requirements, inserted predecessors, supporting operations, and remaining gaps. To relate bridging to outcomes, it must show the provider, endpoints, accepted path, and resolved data bindings. To relate recovery to outcomes, it must distinguish failed attempts, repaired arguments, changed plans, and verified reuse.

The illustrative case provides a concrete explanation of these relationships. It is not a sampled failure distribution or a measured estimate of component utilization. Repeated matched runs and per-task alignment are needed before reporting uncertainty, significance, or causal component effects.

\subsection{Complete Component Results}
\label{app:component-details}
Tables~\ref{tab:ablation-full} and~\ref{tab:ablation-process} retain the complete earlier configuration comparisons. The subset covers all ten domains, with 22 single-step, 20 two- to three-step, and 15 longer records. Its multistep fraction is 61.40\%, so it differs from the main evaluation population. The recorded A1, A2, and A5 intervention limitations remain applicable to these numbers.
\begin{table*}[t]
  \caption{Complete earlier configuration comparisons on the fixed 57-record subset. A1 is a partial control, A2 did not activate the intended ablation, and A5 changes only part of tool ranking. $\Delta$FAS is the full method's mean FAS minus the corresponding configuration's mean FAS, in percentage points; positive values indicate lower performance than the full method.}
  \label{tab:ablation-full}
  \centering
  \begin{tabularx}{\textwidth}{@{}Xrrr@{}}
    \toprule
    Configuration & Mean FAS & Multistep FAS & $\Delta$FAS (pp) \\
    \midrule
    Full method & 0.6387 & 0.5333 & --- \\
    A1: Direct tool-level planning & 0.4748 & 0.4095 & 16.39 \\
    A2: No predecessor expansion & 0.4760 & 0.4095 & 16.27 \\
    A3: No independent completion & 0.4624 & 0.3802 & 17.63 \\
    A4: No directed bridging & 0.6160 & 0.4913 & 2.27 \\
    A5: No tool-selection optimization & 0.6415 & 0.5184 & $-0.28$ \\
    A6: No execution recovery & 0.5972 & 0.4613 & 4.15 \\
    \bottomrule
  \end{tabularx}
\end{table*}

\begin{table*}[t]
\caption{Ablation process metrics. Multistep scores cover the 35 subset records with at least two reference calls. Each column averages per-task scores over its corresponding records. Full uses the supplementary 57-task experiment; A1--A6 use per-task records from the formal 57-task worksheet.}
\label{tab:ablation-process}
\centering
\begin{tabularx}{\textwidth}{@{}Xrrrr@{}}
\toprule
Configuration & Mean TSF1 & Mean TOS & Multistep TSF1 & Multistep TOS \\
\midrule
Full method & 0.9249 & 0.9392 & 0.8862 & 0.8952 \\
A1: Direct tool-level planning & 0.8599 & 0.8778 & 0.8147 & 0.8295 \\
A2: No predecessor expansion & 0.8694 & 0.8958 & 0.8302 & 0.8588 \\
A3: No independent completion & 0.8732 & 0.8865 & 0.8364 & 0.8438 \\
A4: No directed bridging & 0.8603 & 0.8997 & 0.8249 & 0.8652 \\
A5: No tool-selection optimization & 0.8503 & 0.8851 & 0.8133 & 0.8414 \\
A6: No execution recovery & 0.8208 & 0.8588 & 0.7652 & 0.7985 \\
\bottomrule
\end{tabularx}
\end{table*}


Full's overall TSF1/TOS are 0.9249/\allowbreak 0.9392, and its multistep scores are 0.8862/\allowbreak 0.8952. A3 obtains multistep TSF1/TOS of 0.8364/\allowbreak 0.8438, while A4 obtains 0.8249/\allowbreak 0.8652. The higher order score but lower coverage score of A4 illustrates why the two measures should be interpreted separately. These are configuration-level differences rather than an additive decomposition of component effects.

A5 changes only part of tool ranking in the reported run. Its overall FAS rises from Full's 0.6387 to 0.6415, while multistep FAS falls from 0.5333 to 0.5184. Its multistep TSF1/TOS are 0.8133/\allowbreak 0.8414. The directions therefore differ across the answer-score populations.

A6 disables execution recovery and obtains overall/multistep FAS of 0.5972/\allowbreak 0.4613. Its multistep TSF1/TOS are 0.7652/\allowbreak 0.7985, decreases of 12.10/\allowbreak 9.67 percentage points from Full. Attributing these differences to recovered failures requires matched execution traces. Recovery is part of the shared execution framework, so this comparison is supplementary to the main planning interventions.
